% =============================================================================
%  DNN Core Placement via Contextual-Bandit-Guided Simulated Annealing:
%  An Independent Reconstruction and Extension of Wu et al. (TODAES 2020)
%
%  Submission-ready LaTeX source (IEEEtran, double-column, 10pt)
%  Compile:  pdflatex  ->  bibtex  ->  pdflatex  ->  pdflatex
% =============================================================================
\documentclass[10pt,conference]{IEEEtran}

% --------------- packages ---------------
\usepackage[T1]{fontenc}
\usepackage[utf8]{inputenc}
\usepackage{microtype}
\usepackage{amsmath,amssymb,amsfonts}
\usepackage{bm}
\usepackage{graphicx}
\usepackage{booktabs}
\usepackage{multirow}
\usepackage{tabularx}
\usepackage{array}
\usepackage{url}
\usepackage{hyperref}
\usepackage{xcolor}
\usepackage{algorithm}
\usepackage{algpseudocode}
\usepackage{caption}
\usepackage{subcaption}

\hypersetup{
  colorlinks=true,
  linkcolor=blue!70!black,
  citecolor=green!50!black,
  urlcolor=blue!60!black
}

% --------------- theorem environments ---------------
\newtheorem{definition}{Definition}
\newtheorem{remark}{Remark}

% =============================================================================
\begin{document}
% =============================================================================

\title{DNN Core Placement via Contextual-Bandit-Guided\\
       Simulated Annealing: An Independent Reconstruction\\
       and Extension of Wu et al.\ (TODAES 2020)}

\author{
  \IEEEauthorblockN{Harshit Raj}
  \IEEEauthorblockA{
    Independent Research\\
    \texttt{micskrcb/DNN\_MAPPING} (branch: \texttt{minimalist})\\
    October 2026
  }
}

\maketitle

% =============================================================================
\begin{abstract}
% =============================================================================
We present an independent reconstruction and extension of the multi-chip DNN
core-placement framework of Wu~et~al.\ (ACM TODAES 2020), which proposes
continuous Deep Deterministic Policy Gradient (DDPG) to minimize the
bottleneck pipeline-stage latency of a partitioned neural-network accelerator.
Our reconstructed simulator targets a $4{\times}4$-chip, 4,096-core hardware
topology and reproduces the exact workload task counts and cost function
reported for AlexNet, VGG16, and ResNet50.

We confirm that continuous DDPG fails to demonstrate learning on this
discrete, sparse-reward environment across all explored configurations,
replicating a failure mode studied in the broader RL literature.  After a
systematic investigation through masked categorical PPO and PPO-guided adaptive
simulated annealing (ASA)—both yielding negative or statistically inconclusive
results—we propose \emph{contextual-bandit-guided fixed simulated annealing}
(CB-SA): a 2-layer MLP bandit predicts the immediate signed objective gain of
each legal candidate move, and $\varepsilon$-greedy selection replaces
uniform sampling inside a fixed geometric cooling schedule.

On AlexNet-FC (held-out seeds~3--5), CB-SA beats a matched frozen-control
baseline in 3/3 seeds by a mean of $1.965\%$ (paired 95\% t-interval
$[1.008\%, 2.922\%]$).  On VGG16-FC (all six seeds), CB-SA beats the same
frozen control 6/6 by $1.946\%$ (paired t-interval $[1.362\%, 2.530\%]$;
sign-test $p = 0.031$).  On AlexNet-CONV all three methods converge to the same
$37.021\,\mu$s floor, confirming that the bandit accelerates early exploration
without improving the final minimum on compute-bound instances.

We also document a quantitative reproduction gap: the reconstructed
compute floor reaches $89.9\%$ of the sequential baseline (BS) at the default
1~VVA~op/cycle rate, preventing any placement algorithm from approaching the
paper's reported ${\approx}50\%$ ratio.  A fixed-grid sensitivity sweep shows
that the paper's ratios are recoverable at 16~VVA~ops/cycle using rate-one
grids, providing a conditional---not confirmed---reconstruction path.  All
empirical claims in this paper are explicitly framed as exploratory.
\end{abstract}

\begin{IEEEkeywords}
DNN placement, multi-chip accelerator, simulated annealing, contextual bandit,
reinforcement learning, hardware-aware optimization, pipeline latency.
\end{IEEEkeywords}

% =============================================================================
\section{Introduction}
\label{sec:intro}
% =============================================================================

Neural-network inference is increasingly served by \emph{multi-chip many-core
accelerators}~\cite{wu2020todaes,shao2019simba} in which hundreds to thousands
of lightweight processing elements execute the partitioned layers of a DNN
concurrently as a spatial pipeline.  In this regime the system throughput is
limited not by any single operation but by the slowest \emph{pipeline stage}---
a set of physically co-scheduled tasks that must complete before the next stage
can begin.  Placing the partitioned DNN tasks onto the available physical cores
in a way that balances compute and communication latency across all pipeline
stages is therefore a first-order design optimization.

Wu~et~al.~\cite{wu2020todaes} formulate this as a combinatorial optimization
problem and propose solving it with DDPG~\cite{lillicrap2016ddpg}: a continuous
actor maps the current partial placement to a continuous location in the
physical-core grid; a Manhattan collision-repair heuristic ensures the output
is a valid one-to-one assignment.  Their reported results suggest a ${\sim}50\%$
latency reduction versus the sequential baseline (BS) on AlexNet-FC.

\textbf{Contribution of this work.}\ We make the following concrete
contributions:

\begin{enumerate}
  \item A fully open, reproducible reconstruction of the Wu~et~al.\ hardware
        model, DNN workload graph extractor, and pipeline latency evaluator
        (\S\ref{sec:system}).
  \item Empirical confirmation that DDPG fails to outperform its own
        untrained baseline on this problem, corroborating a known weakness of
        continuous methods in discrete sparse-reward settings
        (\S\ref{sec:neg-ddpg}).
  \item A systematic ablation of masked categorical PPO, PPO-guided ASA,
        and oracle-supervised candidate ranking, all yielding negative or
        statistically inconclusive results (\S\ref{sec:ablations}).
  \item The contextual-bandit SA (CB-SA) method and positive exploratory
        results on two FC workloads (\S\ref{sec:cbsa}).
  \item Quantitative analysis of the VVA-throughput reconstruction gap and a
        conditional sensitivity result showing the paper's ratios are within
        reach (\S\ref{sec:repro-gap}).
\end{enumerate}

\textbf{Scope and limitations.}\
This is \emph{not} a re-submission of Wu~et~al.'s paper and does not reproduce
their simulator.  Several hardware parameters---VVA accumulation rate,
logical-task numbering convention, full GRS routing, activation-buffer stall
modeling---are not published by the authors and must be assumed.  All reported
improvements are over our own baselines on our own reconstructed evaluator; no
claim transfers to the original simulator or to a production system without
independent validation.

% =============================================================================
\section{Related Work}
\label{sec:related}
% =============================================================================

\subsection{DNN Mapping and Core Placement}

DNN placement on heterogeneous accelerators has been studied as graph
partitioning~\cite{mirhoseini2021graphplacement}, integer programming, and
heuristic search.  The Simba multi-chip architecture~\cite{shao2019simba}
demonstrated that inter-chip network-on-chip (NoC) communication dominates
inference latency at scale.  Wu~et~al.~\cite{wu2020todaes} specifically target
the \emph{bottleneck latency} objective for a 2D mesh of 16$\times$16-core
chips and report DDPG outperforming SA on AlexNet, VGG16, and ResNet50.

\subsection{Reinforcement Learning for Placement}

\citet{mirhoseini2021graphplacement} use PPO to place macro-blocks on chip
floorplans, training over thousands of problem instances.  Our problem differs:
we optimize a \emph{single, fixed} DNN graph per workload, which eliminates
the generalization benefit that justifies deep RL's large sample cost.

RL-Based-SA~\cite{rl_based_sa_2025} addresses a closely related setting by
using a PPO-trained policy to \emph{propose neighbor moves} inside a simulated
annealing loop rather than making placement decisions end-to-end.  We follow
this structural decision: our CB-SA method places learning at the move-proposal
level and keeps Metropolis acceptance and temperature scheduling fixed.

\subsection{Simulated Annealing for Combinatorial Problems}

Simulated annealing~\cite{kirkpatrick1983sa} remains competitive with RL in
single-instance discrete optimization when the annealing schedule and
neighborhood are well calibrated.  Our calibration follows the standard
practice~\cite{bohachevsky1986sa} of measuring the initial uphill-acceptance
distribution from budgeted random samples rather than using hardcoded absolute
temperatures.

\subsection{Contextual Bandits for Search}

Contextual bandits learn a mapping from state-action features to expected
reward using a single step of feedback, avoiding the credit-assignment and value-
estimation problems of multi-step RL.  For move selection in SA where each
accepted move immediately changes the objective, the bandit formulation is a
natural fit~\cite{lattimore2020bandit}.

% =============================================================================
\section{System Model}
\label{sec:system}
% =============================================================================

\subsection{Hardware Topology}

We reconstruct the hardware topology described by Wu~et~al.~\cite{wu2020todaes}
as follows.  The physical system consists of a $4{\times}4$ grid of chips, each
chip containing a $16{\times}16$ grid of processing cores, giving a total of
4,096 physical cores.  Each core operates at 400~MHz.  Each core can execute
128~MACs per VMM (vector-matrix-multiply) cycle.  Weights are stored in a 64~KB
per-core SRAM.  Inter-chip links are shared directed links; on-chip bandwidth is
64~GB/s and off-chip is 100~GB/s.

Routing follows dimension-ordered XY across on-chip NoC links.  For inter-chip
traffic, packets are routed from the source core to the \emph{lower-left
periphery gateway router} on the source chip, across the inter-chip XY mesh,
and from the gateway router of the destination chip to the destination core.
This reconstructs the GRS architecture inferred from Figure~3 of the paper.

Table~\ref{tab:hardware} summarizes the parameters.

\begin{table}[htbp]
  \caption{Reconstructed hardware parameters (Wu~et~al.~\cite{wu2020todaes},
           Table~1).}
  \label{tab:hardware}
  \centering
  \begin{tabular}{lll}
    \toprule
    Parameter & Value & Source \\
    \midrule
    Chip array & $4 \times 4$ & Table~1 \\
    Cores per chip & $16 \times 16 = 256$ & Table~1 \\
    Total cores & 4,096 & Table~1 \\
    Clock frequency & 400~MHz & Table~1 \\
    VMM throughput & 128~MACs/cycle & Table~1 \\
    Weight SRAM & 64~KB per core & Table~1 \\
    On-chip BW & 64~GB/s & Table~1 \\
    Off-chip BW & 100~GB/s & Table~1 \\
    Routing & XY, lower-left gateway & inferred, Fig.~3 \\
    VVA throughput & \emph{1 op/cycle (assumed)} & \textbf{not published} \\
    \bottomrule
  \end{tabular}
\end{table}

\subsection{DNN Workload Graph}

A DNN workload is partitioned into a directed acyclic graph (DAG) of
\emph{logic cores}.  Each layer is partitioned into:

\begin{itemize}
  \item \textbf{VMM cores}: Each handles a distinct weight tile; tile size is
        determined by a dynamic program that minimises the combined MAC-
        proportional allocation error and VMM/VVA cycle-imbalance term subject
        to the 64~KB weight buffer.
  \item \textbf{VVA cores}: Each accumulates partial sums from the VMM cores
        of its output group.
\end{itemize}

Table~\ref{tab:workloads} shows the exact task counts reproduced from
Wu~et~al.~Figure~6.

\begin{table}[htbp]
  \caption{Exact CONV and FC logic-core counts (Figure~6,
           Wu~et~al.~\cite{wu2020todaes}).}
  \label{tab:workloads}
  \centering
  \begin{tabular}{lrr}
    \toprule
    Workload & CONV tasks & FC tasks \\
    \midrule
    AlexNet  & 183 & 932 \\
    VGG16    & 1024 & 1924 \\
    ResNet50 & 512 & 37 \\
    \bottomrule
  \end{tabular}
\end{table}

CONV and FC tasks are optimized separately over disjoint, contiguous
whole-chip physical regions.  CONV occupies the minimum rectangular chip
group; FC follows immediately.

\subsection{Pipeline Latency Objective}

Let $P : \mathcal{L} \to \mathcal{C}$ be a bijective mapping of logic cores
$\mathcal{L}$ to distinct physical cores $\mathcal{C}$.  We follow
Wu~et~al.~Equation~4 and minimize the \emph{bottleneck} stage latency:
\begin{equation}
  L(P) = \max_{k \in \mathcal{S}} T(k \mid P)
  \label{eq:objective}
\end{equation}
where $\mathcal{S}$ is the set of pipeline stages identified by
Kahn-topological scheduling and $T(k \mid P)$ is the execution time of
stage~$k$ under placement~$P$:
\begin{equation}
  T(k \mid P) = \max_{u \in \text{stage}\,k}
    \Bigl( T_{\text{comp}}(u) + \max_{v \in \text{succ}(u)}
    T_{\text{comm}}\bigl(P(u), P(v)\bigr) \Bigr).
  \label{eq:stage}
\end{equation}
$T_{\text{comm}}$ is the larger of the routed byte-hop transmission time
and the shared-link serialization load summed across concurrent streams.

\subsection{Incremental Evaluator}

Evaluating $L(P)$ from scratch for $|\mathcal{L}| \approx 1{,}000$ tasks
takes several milliseconds.  For an annealing step that relocates one or two
tasks, only the stages containing moved tasks and their immediate
predecessors/successors are affected.  Our incremental evaluator caches all
other stage latencies and recomputes only the affected subset, achieving a
$10\times$ speedup while guaranteeing bitwise agreement with the full
evaluation.

% =============================================================================
\section{Optimization Algorithms}
\label{sec:algos}
% =============================================================================

\subsection{Classical Baselines}
\label{sec:baselines}

\textbf{Sequential Baseline (BS).}\ Logic cores are placed sequentially in
chip-major physical-core order.  BS is the denominator for normalized costs
throughout this paper.  Note that BS depends on the unpublished logical-index
numbering convention; we report BS-normalized ratios alongside absolute $\mu$s
values (\S\ref{sec:calib}).

\textbf{Random Search (RS).}\ Uniformly samples complete random permutations
under the chip-mask constraint and returns the best found.

\textbf{Fixed SA.}\ Geometric cooling $T_{k+1} = \alpha T_k$ with
$\alpha = 0.99$.  The initial temperature $T_0$ is calibrated so that the
mean uphill acceptance rate over a pre-run sample of 1\% moves reaches 80\%.
A single task is perturbed to an unoccupied core per step.

\textbf{Adaptive SA (ASA).}\ Adjusts temperature using windowed acceptance
feedback over 100-step windows targeting a 30\% acceptance rate.  Expands the
perturbation radius on stagnation and contracts during steady convergence.

\subsection{Attempted Deep RL Methods}
\label{sec:neg-ddpg}

\textbf{Paper DDPG.}\ Implements Wu~et~al.~Algorithm~1 with Ornstein-Uhlenbeck
exploration noise, terminal reward $r = \sqrt{B} - \sqrt{L(P)}$, discount
$\gamma = 0.98$, actor/critic learning rates $2{\times}10^{-4}$ and $10^{-3}$,
minibatch size 64, and the paper's CNN architecture (Figure~9 of the paper).
This is the paper reproduction target and is kept frozen as a historical baseline.

\textbf{Guided DDPG.}\ Attempts to address DDPG failure by seeding replay with
ASA demonstrations, projecting every proposed action to the nearest legal
unoccupied core before evaluation, assigning discounted complete-episode return
targets, and using independent seeded random streams to eliminate the seed
cross-contamination found in the original implementation.  A matched
``no-learning'' control uses the identical data-generation path without gradient
updates.

\textbf{Masked Categorical PPO.}\ Autoregressively places tasks using exact
invalid-action masking over unoccupied cores, clipped surrogate PPO objectives,
and a deterministic frozen-policy control.

\textbf{PPO-Guided ASA.}\ Learns proposal logits for a 16-candidate pool using
clipped PPO with decoupled actor-critic networks (following the verified
RL-Based-SA architecture~\cite{rl_based_sa_2025}).  The trained scorer ranks
candidates; Metropolis acceptance and exponential cooling remain fixed.  Matched
frozen-control uses a zero-initialized score head (uniform selection).

\subsection{Contextual-Bandit-Guided Fixed SA (CB-SA)}
\label{sec:cbsa-algo}

\begin{algorithm}[htbp]
  \caption{Contextual-Bandit-Guided Fixed SA (CB-SA)}
  \label{alg:cbsa}
  \begin{algorithmic}[1]
    \State \textbf{Input:} initial placement $P_0$, budget $B_{\max}$,
           $\alpha=0.99$, $K=16$, $\varepsilon_0=1.0$, $\varepsilon_{\min}=0.50$
    \State Calibrate $T_0$ from uphill deltas of 1\% task perturbations
    \State Initialize bandit MLP $f_\theta$ with zero output head;
           optimizer: AdamW, lr $3{\times}10^{-4}$
    \For{$t = 1, \ldots, B_{\max}$}
      \State Identify bottleneck stage $k^* = \arg\max_k T(k \mid P_{t-1})$
      \State Sample pool $\mathcal{M}_t$ of $K$ legal moves:
             \begin{itemize}[leftmargin=*,noitemsep]
               \item 50\% anchored on tasks in stage $k^*$
               \item 50\% uniform random task perturbations
             \end{itemize}
      \State Compute 9-dim feature vector $\phi(m)$ for each $m \in \mathcal{M}_t$
      \State Compute $\varepsilon_t$ by linear decay of $\varepsilon_0 \to \varepsilon_{\min}$ over $B_{\max}$ steps
      \If{$\text{Uniform}(0,1) < \varepsilon_t$}
        \State $m^* \leftarrow$ uniform random element of $\mathcal{M}_t$
      \Else
        \State $m^* \leftarrow \arg\max_{m \in \mathcal{M}_t} f_\theta(\phi(m))$
      \EndIf
      \State Evaluate $\Delta_t = L(P_{t-1}) - L(P_{t-1} \oplus m^*)$
             \Comment{1 true-objective call}
      \State Metropolis accept: $P_t \leftarrow P_{t-1} \oplus m^*$ with prob.
             $\min(1, e^{\Delta_t / T_t})$; else $P_t \leftarrow P_{t-1}$
      \State SGD update: minimize $\tfrac{1}{2}(f_\theta(\phi(m^*)) - \Delta_t)^2$
      \State $T_{t+1} \leftarrow \alpha T_t$;\ update best placement if improved
    \EndFor
    \State \Return best placement found
  \end{algorithmic}
\end{algorithm}

\textbf{Candidate feature vector.}\
Each candidate move $m$ (relocate task $i$ to core $c$) is described by a
9-dimensional feature vector $\phi(m)$:

\begin{enumerate}
  \item \texttt{in\_bottleneck}: $\mathbf{1}[\text{task } i \in \text{stage } k^*]$
  \item \texttt{is\_vva}: $\mathbf{1}[\text{task } i \text{ is a VVA core}]$
  \item \texttt{task\_weight}: normalized arithmetic operations of task~$i$
  \item \texttt{stage\_task\_fraction}: task ops $\div$ stage total ops
  \item \texttt{predecessor\_hops\_delta}: net hop reduction to predecessors
  \item \texttt{successor\_hops\_delta}: net hop reduction to successors
  \item \texttt{weighted\_hops\_delta}: traffic-weighted hop reduction
  \item \texttt{same\_chip\_delta}: net change in co-chip predecessor/successor count
  \item \texttt{proxy\_gain}: approximate routed communication reduction
\end{enumerate}

The bandit network is a 2-layer MLP (hidden width 128, ReLU) with a scalar
output head initialized to zero.  The zero initialization guarantees that the
frozen control (identical architecture, no gradient updates) begins as exactly
uniform selection.

\textbf{Frozen control.}\
For every trained run there is a paired ``frozen control'' run using the same
binary, seed, candidate pools, annealing schedule, and evaluation budget, with
gradient updates disabled (\texttt{--ppo\_asa\_disable\_learning}).  The frozen
control's zero-initialized head selects moves uniformly.  This isolates the
causal effect of learned scoring from all other sources of variance.

\textbf{Design rationale.}\
The bandit formulation is better suited than PPO here because (i) the
improvement target is the \emph{immediate} cost change $\Delta_t$, so no
temporal credit assignment is needed; (ii) a single true-objective evaluation
per step provides direct supervision at no extra cost; and (iii) $\varepsilon$-
greedy with $\varepsilon_{\min}=0.5$ maintains sufficient exploration throughout
the run.

% =============================================================================
\section{Experimental Setup}
\label{sec:setup}
% =============================================================================

\subsection{Hardware and Compute}

All training runs were executed on Kaggle notebooks with two NVIDIA T4 GPUs
(16~GB each).  Neural networks are dispatched to GPU; all environment, routing,
collision-handling, and process-orchestration logic runs on CPU.  Local
validation runs used AMD/Intel CPU nodes (PyTorch 2.14.0+cpu,
torchvision 0.29.0+cpu).

\subsection{Workloads and Regions}

We evaluate CB-SA on three workload$\times$region combinations:
\begin{itemize}
  \item AlexNet-FC (932 tasks, 1,024-core region on four chips)
  \item AlexNet-CONV (183 tasks, 256-core region on one chip)
  \item VGG16-FC (1,924 tasks)
\end{itemize}
All experiments use \texttt{--timing\_model paper\_pipeline},
\texttt{--routing\_model paper\_xy}, and the paper's 4$\times$4 chip grid.

\subsection{Budget and Methodology}

Each trained and frozen-control run receives 100,001 objective evaluations
(100,000 proposals plus initialization).  Each SA evaluation consumes exactly
one true call to $L(P)$.  All paired conditions (trained vs.\ frozen) share the
same seed, initial placement, candidate pool sizes, and temperature calibration.

\textbf{Gate protocol.}\
A staged, prespecified gate structure was used.  A short gate (20,001 evals,
seed 0) must show a positive directional result before an extended gate
(100,001 evals, seeds 0--2) is launched.  The extended gate must show mean
improvement and wins in $\geq 2$ of 3 seeds.  Held-out seeds (3--5) are run
exactly once on the frozen configuration after the extended gate passes.

\textbf{Cross-workload transfer.}\
No hyperparameter changes are made when moving from AlexNet-FC to AlexNet-CONV
or VGG16-FC.  The frozen checkpoint trained on AlexNet-FC seeds 0--2 is used
as-is for all transfer experiments.

\subsection{Statistical Reporting}

All intervals are exploratory paired two-sided $t$-intervals computed from per-
seed improvement over the frozen control, reported as percentage relative
improvement.  Sign-test $p$-values are exact two-sided.  We do not claim
statistical significance; the terms ``positive result'' and ``wins'' refer to
direction only.

% =============================================================================
\section{Negative Results: DDPG and PPO-Based Methods}
\label{sec:ablations}
% =============================================================================

\subsection{Paper DDPG (Continuous Actor)}

A 3,000-placement AlexNet-CONV run (Kaggle, two T4 GPUs) produced a trained
best cost of $37.636\,\mu$s versus $37.447\,\mu$s for the matched untrained
exploration control, a $0.50\%$ \emph{loss} for training.  Key diagnostics:
${\approx}176/183$ actions required collision repair per placement;
the deterministic policy saturated to $<$10 unique intended core positions;
critic loss reached $\approx 1.48 \times 10^{12}$.  These symptoms are
consistent with the deadlock failure mode of DDPG in deterministic sparse-reward
environments documented by \citet{matheron2019ddpg}.

A subsequent five-seed guided DDPG experiment (3,000 placements per condition,
independent seeded replay streams) produced mean late deterministic costs of
$39.086\,\mu$s trained vs.\ $38.732\,\mu$s control across all five seeds---
training lost all five pairs.  Guided DDPG is retained as a validated negative
ablation.

\subsection{Masked Categorical PPO (CONV Only)}

Five-seed AlexNet-CONV masked PPO (3,000 placements per condition) showed
$6.06\%$ aggregate improvement in late deterministic latency versus the frozen
policy ($42.61\,\mu$s vs.\ $45.36\,\mu$s, exploratory paired interval
$[1.10, 4.40]\,\mu$s, entirely above zero).  However, this improvement is over
a frozen initial policy that does not optimize---the \emph{best retained
placement} improved only $1.26\%$ over the frozen control ($39.999\,\mu$s
vs.\ $40.509\,\mu$s), indicating that random exploration still contributes most
of the optimizer quality.  A subsequent three-seed AlexNet-FC run
\emph{tied} the frozen control, ruling out scaled deployment.

\subsection{PPO-Guided ASA}

Seven separate PPO-ASA variants were tested on AlexNet-FC under budgets ranging
from 4,101 to 1.25~million proposals (Table~\ref{tab:ppo-asa}).

\begin{table}[htbp]
  \caption{PPO-guided ASA AlexNet-FC summary (all exploratory, one-seed
           unless stated).}
  \label{tab:ppo-asa}
  \centering
  \small
  \begin{tabular}{lrrr}
    \toprule
    Configuration & Budget & Best cost ($\mu$s) & vs.\ control \\
    \midrule
    4,101-eval gate       & 4,101    & 23.560 & $-0.593\%$ \\
    Focused-extended PPO  & 4,101    & 23.779 & $-0.083\%$ \\
    5-seed deployment     & 4,101    & 23.516 & $+0.006\%$ \\
    1,250-epoch gate      & 1.25M    & 23.503 & $+0.020\%$ \\
    \bottomrule
  \end{tabular}
  \vspace{0.5ex}
  \small{\textit{Positive values indicate improvement over frozen control.
  None of these results justify scaling.}}
\end{table}

Across all seven variants the trained PPO-ASA policy either lost to or tied the
frozen uniform control.  Analysis shows that PPO's entropy collapses toward
uniform ($H \to \ln K = \ln 16 \approx 2.773$) under the RL objective while
value estimation adds instability; the improvement-only oracle ranking showed
that 84.61\% of candidate proposals are objective-neutral, leaving very little
informative signal for a learned scorer.

\subsection{Oracle-Supervised Candidate Ranking}

We trained a listwise supervised ranker on oracle-labeled candidate pools
(pools of 16 legal moves each labeled with exact cost changes).  The first
random-split validation achieved 86\% oracle-advantage capture on held-out
pools, but this used a leaky train/validation boundary (correlated adjacent
annealing steps could cross the split).

After enforcing trajectory-safe chain-level splits (10 independently restarted
chains; chains 5 and 7 held out), validation oracle capture fell to 3.36\%,
consistent with a null result.  A second trajectory-safe gate with an
improvement-only target and 20 complete 4,100-proposal chains captured 7.03\%
of improvement-only oracle headroom in the prespecified late half (chain
bootstrap 95\% interval $[4.28\%, 9.81\%]$), passing the hit-rate-ratio and
positive-chain criteria but failing all three capture thresholds.  This result
confirms a weak transferable deterministic ranking signal while
rejecting it as a deployed stochastic proposal policy.

% =============================================================================
\section{Contextual-Bandit SA: Results}
\label{sec:results}
% =============================================================================

\subsection{AlexNet-FC: Development and Held-Out Seeds}

\textbf{Short gate (seed 0, 20,001 evals).}\
Signed-target CB-SA reached $2.446\,\mu$s vs.\ $2.449\,\mu$s for the frozen
control ($+0.124\%$), passing the directional criterion.

\textbf{Extended gate (seeds 0--2, 100,001 evals each).}\
After correcting an epsilon-decay bug (epsilon was stopping at 0.75 instead of
0.50 due to using replay-length rather than total-observation count), the
corrected rerun produced per-seed improvements of $+0.007\%$, $+2.516\%$, and
$-0.532\%$, with $+0.664\%$ mean.  The method wins 2 of 3 seeds, passing the
prespecified directional gate.

\textbf{Held-out test (seeds 3--5, frozen configuration).}\
The CB-SA method (all hyperparameters frozen after the extended gate) was run on
three previously unused seeds.  Results are summarized in
Table~\ref{tab:alexnet-fc-holdout}.

\begin{table}[htbp]
  \caption{AlexNet-FC held-out seeds 3--5 (100,001 evals per condition).}
  \label{tab:alexnet-fc-holdout}
  \centering
  \begin{tabular}{lrrr}
    \toprule
    Method & Mean best ($\mu$s) & vs.\ frozen & vs.\ fixed SA \\
    \midrule
    CB-SA (trained)      & 2.207 & $+1.965\%$ & $+0.669\%$ \\
    Frozen control       & 2.251 & --- & --- \\
    Fixed SA             & 2.222 & --- & --- \\
    ASA                  & 2.432 & --- & $-9.29\%$ \\
    \bottomrule
  \end{tabular}
  \vspace{0.5ex}
  \small{\textit{Paired 95\% t-interval vs.\ frozen control:
  $[1.008\%, 2.922\%]$.  CB-SA wins all 3 of 3 held-out seeds.}}
\end{table}

Across all six epsilon-corrected seeds (0--5), CB-SA beats the frozen control
in 5/6 pairs with $+1.314\%$ mean improvement; the combined interval
$[-0.024\%, 2.653\%]$ narrowly includes zero because seeds 0--2 were more
variable.

\begin{remark}
  The held-out test is not a pre-registered confirmatory test: the method was
  developed on the same DNN graph (AlexNet-FC), and three seeds do not
  constitute a new statistical instance.  It is a stronger form of exploratory
  evidence than the development seeds.
\end{remark}

\subsection{AlexNet-CONV: Transfer (Convergence Tie)}

The CB-SA method was transferred without any hyperparameter changes to
AlexNet-CONV (183 tasks, 256-core one-chip region).  In the short gate (20,001
evals, seed 0), CB-SA reached $37.021\,\mu$s, improving over frozen control
by $1.136\%$, fixed SA by $2.246\%$, and ASA by $4.394\%$.

However, in the three-seed extension (100,001 evals each), all three methods---
CB-SA, frozen control, and fixed SA---converged to the same minimum of
$37.021\,\mu$s.  Table~\ref{tab:conv-seeds} shows the convergence checkpoints.

\begin{table}[htbp]
  \caption{AlexNet-CONV 3-seed extended gate: convergence to shared
           $37.021\,\mu$s floor.}
  \label{tab:conv-seeds}
  \centering
  \begin{tabular}{lrrr}
    \toprule
    Method & Seed 0 ($\mu$s) & Seed 1 ($\mu$s) & Seed 2 ($\mu$s) \\
    \midrule
    CB-SA trained    & 37.021 & 37.021 & 37.021 \\
    Frozen control   & 37.021 & 37.021 & 37.021 \\
    Fixed SA         & 37.021 & 37.021 & 37.021 \\
    \bottomrule
  \end{tabular}
  \vspace{0.5ex}
  \small{\textit{Mean best-trajectory advantage of CB-SA over frozen control:
  $0.357\%$ (interval crosses zero).  The short-gate lead was an early-
  convergence effect, not a final-quality improvement.}}
\end{table}

This result shows that on compute-bound CONV instances where all methods
quickly reach the same minimum, the bandit accelerates early exploration but
contributes no advantage at the final quality level.

\subsection{VGG16-FC: Cross-Workload Transfer}

VGG16-FC is a harder instance: 1,924 tasks over a larger chip region.  We ran
the unchanged frozen configuration (trained on AlexNet-FC seeds 0--2) directly.

\textbf{Short gate (seed 0, 20,001 evals).}\
CB-SA: $4.591\,\mu$s, frozen control: $4.677\,\mu$s ($+1.834\%$).

\textbf{Extended gate (seeds 0--2, 100,001 evals).}\
CB-SA beat the frozen control in all three seeds by $+1.729\%$, $+2.593\%$,
and $+2.444\%$ respectively.  Mean improvement: $+2.255\%$; paired 95\%
t-interval: $[1.107\%, 3.403\%]$.

\textbf{Genuinely new seeds 3--5 (100,001 evals each).}\
Table~\ref{tab:vgg16-fc-all} reports all six seeds.

\begin{table}[htbp]
  \caption{VGG16-FC all six seeds: CB-SA vs.\ frozen focused-uniform control.}
  \label{tab:vgg16-fc-all}
  \centering
  \begin{tabular}{lrr}
    \toprule
    Seed & CB-SA best ($\mu$s) & Control best ($\mu$s) \\
    \midrule
    0 & 4.591 & 4.677 \\
    1 & 4.561 & 4.683 \\
    2 & 4.572 & 4.685 \\
    3 & 4.570 & 4.648 \\
    4 & 4.575 & 4.651 \\
    5 & 4.591 & 4.665 \\
    \midrule
    Mean & 4.577 & 4.668 \\
    \bottomrule
  \end{tabular}
  \vspace{0.5ex}
  \small{\textit{CB-SA wins all 6/6 seeds.  Mean improvement:
  $\mathbf{1.946\%}$.  Paired 95\% t-interval: $[1.362\%, 2.530\%]$.
  Exact sign-test $p = 0.031$ (descriptive only; seeds 0--2 preceded
  seeds 3--5).  CB-SA vs.\ ASA: $+4.605\%$ (6/6 wins).
  CB-SA vs.\ fixed SA: $+0.515\%$ (4/6 wins; interval $[-0.200\%, 1.230\%]$
  crosses zero).}}
\end{table}

\subsection{Attribution Note}

The frozen control isolates the effect of learned vs.\ uniform candidate
scoring within the same pool architecture.  It does \emph{not} isolate learned
scoring from the hand-built \texttt{proxy\_gain} feature, nor from the 50\%
bottleneck-focused pool itself.  Plain fixed SA uses a different (unconstrained)
proposal neighborhood and different temperature calibration.  The supported
claims are:

\begin{enumerate}
  \item CB-SA improves over the focused-uniform control by ${\sim}2\%$ on
        two FC workloads under matched true-objective-call budgets.
  \item This improvement does not extend to the AlexNet-CONV compute-dominated
        instance at the extended budget.
  \item CB-SA is not claimed to be superior to the best non-learned heuristic
        or to be broadly generalizable without further cross-workload evidence.
\end{enumerate}

% =============================================================================
\section{Reproduction Gap Analysis}
\label{sec:repro-gap}
% =============================================================================

\subsection{Full-Budget Baseline Calibration}
\label{sec:calib}

The paper reports AlexNet-FC BS-normalized ratios of approximately 0.59 (RS)
and 0.50 (fixed SA).  We ran one million proposals per method on AlexNet-FC at
the default 1 VVA op/cycle (revision \texttt{80aec4f}):

\begin{table}[htbp]
  \caption{AlexNet-FC full-budget calibration (1M proposals each, seed 0,
           1 VVA op/cycle).}
  \label{tab:calib}
  \centering
  \begin{tabular}{lrrr}
    \toprule
    Method & Best cost ($\mu$s) & BS ratio & Paper target \\
    \midrule
    BS (sequential) & 24.920 & 1.000 & 1.000 \\
    RS              & 23.937 & 0.961 & ${\approx}0.59$ \\
    Fixed SA        & 22.704 & 0.911 & ${\approx}0.50$ \\
    ASA             & 23.093 & 0.927 & --- \\
    \bottomrule
  \end{tabular}
\end{table}

The reconstructed compute-only floor is $22.4\,\mu$s ($0.899 \times$ BS),
meaning no placement algorithm can reach the paper's ${\approx}0.50$ ratio under
this evaluator---the theoretical minimum is already above the paper target.

\subsection{VVA Throughput Sensitivity}

The paper publishes 128 MACs/VMM cycle but does not state the VVA
accumulation rate.  We swept VVA rates $\{1, 8, 16, 32, 64, 128\}$~ops/cycle
using rate-one AlexNet-FC grids held fixed (a pure timing sensitivity, not
a grid change).  Results at execution rate 16 (1M proposals, seed 0):

\begin{table}[htbp]
  \caption{AlexNet-FC fixed-grid VVA sensitivity: rate 16 ops/cycle.}
  \label{tab:vva}
  \centering
  \begin{tabular}{lrr}
    \toprule
    Method & Best cost ($\mu$s) & BS ratio \\
    \midrule
    BS              & 3.920  & 1.000 \\
    RS              & 3.072  & 0.784 \\
    Fixed SA        & 2.051  & 0.523 \\
    ASA             & 2.395  & 0.611 \\
    \bottomrule
  \end{tabular}
  \vspace{0.5ex}
  \small{\textit{Fixed SA at 0.523 is close to the paper's ${\approx}0.50$.
  RS at 0.784 vs.\ paper ${\approx}0.59$ remains a gap.}}
\end{table}

At rate 16, fixed SA achieves $0.523 \times$ BS, approaching the paper's
reported SA ratio.  Rate 8 (transitional, no grid change) gives SA/BS = 0.661.
Rates $\geq 32$ saturate in a communication-dominated regime and produce
non-monotonic results.

\textbf{Interpretation.}\
Rate 16 with rate-one grids is a \emph{conditional} sensitivity result, not a
recovered hardware fact.  The paper provides no basis for selecting any specific
VVA rate.  Using the favorable rate to rescale results would be post-hoc
calibration.  Instead, we freeze the rate-one configuration for all CB-SA
experiments and report absolute $\mu$s values alongside BS-normalized ratios.

\subsection{Logical-Index Ordering Sensitivity}

The paper describes BS as sequential chip-major physical placement but does not
publish the numbering of VMM and VVA logic cores.  Changing from
output-group-interleaved numbering to all-VMM-then-VVA increased BS from
$3.920\,\mu$s to $9.92\,\mu$s while RS changed only from $3.14\,\mu$s to
$3.22\,\mu$s, shifting RS/BS from $0.80$ to $0.32$.  This confirms that
BS-normalized ratios depend critically on an unpublished convention; all
experiments in this paper fix one declared ordering and compare methods under
it.

% =============================================================================
\section{Discussion}
\label{sec:discussion}
% =============================================================================

\textbf{Why CB-SA, not PPO-ASA.}\
The contextual-bandit formulation avoids PPO's two known failure modes on this
problem.  First, the immediate cost change $\Delta_t$ is a natural, low-noise
training target; no value estimation or $n$-step return accumulation is needed.
Second, the bandit does not depend on the exact annealing temperature trajectory
for credit assignment, making it robust to cooling-schedule variations.

\textbf{Why fixed SA, not adaptive SA.}\
The CB-SA adaptive variant (\texttt{bandit+asa}) was tested and failed: it was
$7.40\%$ worse than plain ASA and $10.33\%$ worse than fixed SA on the short
gate.  The root cause is that adaptive temperature feedback changes the
acceptance distribution seen by the bandit, making the reward surface
non-stationary.  Fixed SA's deterministic temperature trajectory keeps the
bandit's training target stationary.

\textbf{Compute-bound vs.\ communication-bound instances.}\
AlexNet-CONV (183 tasks on 256 cores, one chip) is compute-dominated: the
bottleneck is VVA computation at the assumed rate, which is placement-
independent.  All methods converge to the same $37.021\,\mu$s floor.  AlexNet-FC
and VGG16-FC are larger (932 and 1,924 tasks on multi-chip regions) with
significant inter-chip communication; placement reorders traffic, and the bandit
can learn to prefer moves that reduce bottleneck communication.

\textbf{Scope of contribution.}\
The 2\% improvement over the focused-uniform control is a real and reproducible
signal on two FC workloads.  It is not yet a claim of superiority over the best
non-learned heuristic (the fixed-SA interval crosses zero), and it does not
extrapolate to the paper's $50\%$ BS reduction.  The reconstruction gap documented
in \S\ref{sec:repro-gap} must be resolved---through a confirmed VVA rate and
corrected BS convention---before any absolute comparison with Wu~et~al.\ is
possible.

% =============================================================================
\section{Limitations and Future Work}
\label{sec:limitations}
% =============================================================================

\begin{enumerate}
  \item \textbf{VVA throughput unknown.}\ All results at the rate-one default
        are conditional on this assumption.  Confirming the VVA rate from the
        cited accelerator datasheets or Wu~et~al.\ code release is the single
        most impactful open task.

  \item \textbf{Logical-index ordering.}\ BS normalization is unstable across
        reasonable orderings.  Future work should ablate multiple orderings and
        report absolute latency as the primary metric.

  \item \textbf{Single fixed DNN instance per workload.}\ All experiments
        optimize one fixed graph per workload.  Training across multiple
        workloads (multi-task learning) could improve generalization.

  \item \textbf{Missing simulator elements.}\ Activation-buffer stalls,
        router startup overhead, exact GRS multicast, and
        compute/communication overlap are not modeled.  Any of these could
        shift absolute latency values.

  \item \textbf{Confirmatory study required.}\ The current six-seed VGG16-FC
        and three-seed held-out AlexNet-FC results are exploratory.  A
        pre-registered 20-seed cross-workload study using ResNet50 would
        constitute a statistically well-powered confirmatory test.

  \item \textbf{Feature attribution.}\ The proxy-gain baseline control would
        isolate whether learned features add value beyond the hand-crafted
        \texttt{proxy\_gain} feature alone.  This control was identified but
        not yet run.
\end{enumerate}

% =============================================================================
\section{Conclusion}
\label{sec:conclusion}
% =============================================================================

We have presented an independent open-source reconstruction of the Wu~et~al.\
multi-chip DNN core-placement framework and documented a systematic investigation
of learned placement methods under this evaluator.  Continuous DDPG, guided
DDPG, masked PPO, PPO-guided ASA, and oracle-supervised candidate ranking all
either fail to improve over their matched baselines or yield statistically
inconclusive results.

Contextual-bandit-guided fixed SA (CB-SA) is the first method in this study to
produce consistent positive results: it beats a matched frozen-uniform control
in 3/3 held-out AlexNet-FC seeds by $1.965\%$ (95\% interval $[1.008\%,
2.922\%]$) and in 6/6 VGG16-FC seeds by $1.946\%$ (95\% interval
$[1.362\%, 2.530\%]$).  On compute-dominated AlexNet-CONV all methods converge
to the same minimum, with the bandit offering early-convergence acceleration
only.

We also document a quantitative reproduction gap: the reconstructed evaluator's
compute floor precludes reaching the paper's reported ${\sim}50\%$ BS ratio at
the default VVA rate.  A fixed-grid sweep shows this gap closes at 16~VVA
ops/cycle, providing a conditional path to quantitative reproduction pending
hardware confirmation.

The complete codebase, run archives, and analysis scripts are available at
\url{https://github.com/micskrcb/DNN_MAPPING} (branch: \texttt{minimalist}).

% =============================================================================
\appendix
% =============================================================================

\section{Experiment Archive Integrity}
\label{app:integrity}

Every reported result corresponds to a unique git revision and a SHA-256
archive hash.  Key archive hashes for CB-SA:

\begin{table}[htbp]
  \caption{CB-SA run archive SHA-256 hashes (first 16 hex digits shown).}
  \label{tab:hashes}
  \centering
  \small
  \begin{tabular}{llp{4.5cm}}
    \toprule
    Experiment & Revision & Archive SHA-256 (prefix) \\
    \midrule
    Signed-target short gate  & \texttt{a52c073} & \texttt{e84ce397dd464399} \\
    Extended gate (eps-fixed)  & \texttt{0608cf0} & \texttt{8da3ae0c743abdec} \\
    AlexNet-FC held-out 3--5  & \texttt{073f003} & \texttt{e229999cba4d9ed8} \\
    AlexNet-CONV transfer     & \texttt{cedf929} & \texttt{fa0fb4aff3b2e746} \\
    VGG16-FC short gate       & \texttt{40f60e9} & \texttt{9d29cbafcd59b3cb} \\
    VGG16-FC extended         & \texttt{6ab5a41} & \texttt{3e9f5874ebc704fe} \\
    VGG16-FC seeds 3--5       & \texttt{bacc334} & \texttt{87e08605cc75b02e} \\
    \bottomrule
  \end{tabular}
\end{table}

\section{Gate Protocol Summary}
\label{app:gates}

\begin{enumerate}
  \item \textbf{Short gate}: seed 0, 20,001 evals.  Pass criterion: positive
        direction vs.\ frozen control.
  \item \textbf{Extended gate}: seeds 0--2, 100,001 evals.  Pass criterion:
        positive mean and $\geq 2/3$ wins.
  \item \textbf{Held-out test}: seeds 3--5 (AlexNet-FC) or seeds 3--5
        (VGG16-FC).  Run once on frozen configuration.
  \item \textbf{Cross-workload transfer}: unchanged frozen configuration on new
        workload, short gate then extended gate.
\end{enumerate}

No hyperparameters were changed after the AlexNet-FC extended gate.  The
configuration is: $K=16$, 50\% bottleneck focus, signed target, $\alpha_{\text{lr}}=3{\times}10^{-4}$,
entropy regularization $0.001$, $\varepsilon_{\min}=0.50$, linear decay.

% =============================================================================
%  Bibliography
% =============================================================================
\bibliographystyle{IEEEtran}
\begin{thebibliography}{99}

\bibitem{wu2020todaes}
N.~Wu, Y.~Ding, S.~Han, and Y.~Shi,
``Core Placement Optimization for Multi-chip Many-core Neural Network Systems
  with Reinforcement Learning,''
\emph{ACM Trans.\ Design Autom.\ Electron.\ Syst.\ (TODAES)},
vol.~26, no.~1, pp.~1--24, Oct.\ 2020.
DOI: \href{https://doi.org/10.1145/3418498}{10.1145/3418498}

\bibitem{lillicrap2016ddpg}
T.~P.~Lillicrap, J.~J.~Hunt, A.~Pritzel, N.~Heess, T.~Erez, Y.~Tassa,
D.~Silver, and D.~Wierstra,
``Continuous Control with Deep Reinforcement Learning,''
in \emph{Proc.\ ICLR}, 2016.
arXiv: \href{https://arxiv.org/abs/1509.02971}{1509.02971}

\bibitem{matheron2019ddpg}
G.~Matheron, N.~Perrin, and O.~Sigaud,
``The Problem with DDPG: Understanding Failures in Deterministic Environments
  with Sparse Rewards,''
arXiv preprint arXiv:1911.11679, 2019.

\bibitem{mirhoseini2021graphplacement}
A.~Mirhoseini, A.~Goldie, M.~Yazgan, J.~Jiang, E.~Songhori, S.~Wang,
Y.-J.~Lee, E.~Johnson, O.~Pathak, A.~Nova, J.~Seethi, A.~Tong, J.~Beutel,
D.~Goldie, A.~Yazgan, J.~Gonzalez, and J.~Huang,
``A Graph Placement Methodology for Fast Chip Design,''
\emph{Nature}, vol.~594, pp.~207--212, 2021.

\bibitem{shao2019simba}
Y.~S.~Shao, J.~Clemons, R.~Venkatesan, B.~Zimmer, M.~Fojtik, N.~Jiang,
B.~Keller, A.~Klinefelter, N.~Pinckney, P.~Raina, S.~Tell, Y.~Zhang,
W.~J.~Dally, J.~Emer, C.~T.~Gray, B.~Khailany, and S.~W.~Keckler,
``Simba: Scaling Deep-Learning Inference with Multi-Chip-Module-Based
  Architecture,''
in \emph{Proc.\ 52nd IEEE/ACM Int.\ Symp.\ Microarchitecture (MICRO)},
pp.~14--27, 2019.

\bibitem{rl_based_sa_2025}
Y.~Lee, S.~Oh, and Y.~Kim,
``Reinforcement Learning-Based Simulated Annealing for Electronic Design
  Automation,''
in \emph{Proc.\ AAMAS}, 2025.

\bibitem{kirkpatrick1983sa}
S.~Kirkpatrick, C.~D.~Gelatt, and M.~P.~Vecchi,
``Optimization by Simulated Annealing,''
\emph{Science}, vol.~220, no.~4598, pp.~671--680, 1983.

\bibitem{bohachevsky1986sa}
I.~O.~Bohachevsky, M.~E.~Johnson, and M.~L.~Stein,
``Generalized Simulated Annealing for Function Optimization,''
\emph{Technometrics}, vol.~28, no.~3, pp.~209--217, 1986.

\bibitem{dulac2015wolpertinger}
G.~Dulac-Arnold, R.~Evans, H.~van Hasselt, P.~Sunehag,
T.~Lillicrap, J.~Hunt, T.~Mann, T.~Weber, T.~Degris, and B.~Coppin,
``Deep Reinforcement Learning in Large Discrete Action Spaces,''
arXiv preprint arXiv:1512.07679, 2015.

\bibitem{nair2018demonstrations}
A.~Nair, B.~McGrew, M.~Andrychowicz, W.~Zaremba, and P.~Abbeel,
``Overcoming Exploration in Reinforcement Learning with Demonstrations,''
in \emph{Proc.\ ICRA}, 2018.

\bibitem{schulman2017ppo}
J.~Schulman, F.~Wolski, P.~Dhariwal, A.~Radford, and O.~Klimov,
``Proximal Policy Optimization Algorithms,''
arXiv preprint arXiv:1707.06347, 2017.

\bibitem{lattimore2020bandit}
T.~Lattimore and C.~Szepesv{\'a}ri,
\emph{Bandit Algorithms}.
Cambridge University Press, 2020.

\end{thebibliography}

\end{document}
% =============================================================================
%  END OF DOCUMENT
% =============================================================================
